For a full distribution, count equally likely permutations, subsets, or gap configurations. For a mean, a sum of indicators often avoids the full law and needs no independence. Variances require joint probabilities, which can be classified by overlap. In circular counting, specify the marking rule and how often each object is counted.

Exchangeability is weaker than independence. Independence created by a Poisson sample size must be proved by factoring the joint law. If the number of terms depends on observations, retain that dependence in a first-step calculation.

A high-dimensional volume can be expressed as a probability for independent uniform samples. The law of large numbers locates a threshold; when the mean cancels and variance is finite and nonzero, the central limit theorem describes fluctuations near it.
\par\smallskip
